% !TeX root = ../paper.tex
\section{Benchmark Evaluation Settings and LIBERO-SMART Protocol}
\label{appendix:smart-benchmark}

\textbf{Benchmark adaptation and evaluation convention.}
Our cross-model comparisons evaluate the downstream capabilities
of trained policies. Models retain their respective pretraining
data, objectives, and training budgets. Except where explicitly
noted, baselines follow the standard training or adaptation
protocol released with each benchmark: they are either SFT-adapted
on the corresponding training data or evaluated using a publicly
released target-domain adapted checkpoint. For evaluations conducted
in this work, we match tasks, initializations, interaction horizons,
and success criteria within each benchmark and evaluation protocol.
Test-time evaluation follows the native benchmark protocol;
LIBERO-SMART uses the protocol specified below. Its evaluation
instances are not used to adapt any model. The scope of the
within-model component ablations is described in
Appendix~\ref{appendix:ablation-details}.

\textbf{LIBERO, LIBERO-Plus, and LIBERO-Pro.} These benchmarks evaluate basic task execution,
perturbation robustness, and task and position generalization, respectively. In our evaluations, we
retain their native goal-based success criteria and interaction horizons, using matched tasks and
initializations across methods within each benchmark.

For each SmartVLA architecture, the same LIBERO-adapted expert is evaluated
on LIBERO, LIBERO-Plus, LIBERO-Pro, and LIBERO-SMART, without further
benchmark-specific adaptation. LIBERO-SMART evaluation instances are not
used in any of these adaptation runs.

\textbf{RoboTwin2.0.} Full adapts on clean and randomized demonstrations and evaluates both scene
types; Clean2Random (C2R) adapts only on clean demonstrations and evaluates clean and randomized
scenes. The two protocols use independently adapted experts.

\textbf{LIBERO-SMART.} We evaluate the three instruction-following capabilities using the
trajectory-level task contracts in Section~\ref{libero-smart}, with matched instances and
interaction budgets across models within each evaluation protocol. The agentic calling protocol and
success aggregation rules are detailed in Appendices~\ref{agentic-calling-protocol}
and~\ref{app:smart-aggregation}.

\textbf{Piper and training.} Piper evaluates instruction-following gains in
real-world bimanual control. SmartVLA uses a separate expert for each target
domain following Section~\ref{steering-centric-policy-training}; model and
execution details are provided in the appendices.

\subsection{Tasks and Initialization Instances}
\label{app:smart-init}

The benchmark comprises three independently scored suites. The task and
initialization counts are summarized in Table~\ref{tab:smart-suite-composition}.

% !TeX root = ../paper.tex
\begin{table}[htbp]
\centering
\setlength{\belowcaptionskip}{\baselineskip}
\caption{LIBERO-SMART task suites and initialization counts.}
\label{tab:smart-suite-composition}
\normalsize
\setlength{\tabcolsep}{3pt}
\renewcommand{\arraystretch}{1.10}
\begin{tabular}{@{}
  >{\raggedright\arraybackslash}p{(\linewidth - 8\tabcolsep) * \real{0.21}}
  >{\raggedleft\arraybackslash}p{(\linewidth - 8\tabcolsep) * \real{0.08}}
  >{\raggedleft\arraybackslash}p{(\linewidth - 8\tabcolsep) * \real{0.12}}
  >{\raggedleft\arraybackslash}p{(\linewidth - 8\tabcolsep) * \real{0.13}}
  >{\raggedright\arraybackslash}p{(\linewidth - 8\tabcolsep) * \real{0.46}}@{}}
\toprule
Suite & Tasks & Init./task & Instances & Core content \\
\midrule
Semantic & 40 & 50 & 2,000 & paraphrases, sub-tasks, atomic motions, constraints \\
Visual Prompt (Vis) & 40 & 50 & 2,000 & object circles, regions, points, trajectories \\
Control & 40 & 50 & 2,000 & target correction and process re-planning \\
\textbf{Total} & \textbf{120} & \textbf{50} & \textbf{6,000} & three capabilities \\
\bottomrule
\end{tabular}
\end{table}
%

Every task has 10 macro layouts and 5 local initialization variants per
layout. For language-reference tasks, the target relation must hold in
the current scene with a unique referent; left/right relations are
resolved from the task-specified view rather than a fixed object
identifier. For visual-prompt tasks, cue position, target entity, and
spatial predicate are updated together for each instance. For Control,
the expected task phase and required held or completed state are
preserved while scene placement and feasible robot pose are varied.

\textbf{Controlled initialization for corrective instruction following.}
For each Control instance, all models begin from the same
intermediate physical state and are evaluated on their response
to the corrective instruction. This fixes the robot configuration,
object arrangement, and required interaction state at the start
of the evaluated continuation. If each model instead executed
the preceding stage independently, differences in grasp pose,
object placement, or incomplete preceding actions could change
the difficulty of the correction. Standardized initialization
controls these preceding-stage effects and permits comparison of
how models execute the updated requirement from the same physical
starting point. Models receive the same corrective instruction
and interaction budget within each evaluation protocol.

Validity checks cover key-object visibility, target reachability,
support and contact relations, and object stability. We exclude
geometric intersections, unreasonable occlusion, or extra obstacles that
change instruction meaning. Tasks requiring a new operation exclude
initial states in which the target is already complete. Holding tasks
retain the required initial state and test whether the policy maintains
it during execution. Initial states are specified reproducibly and
shared across models.

Each Visual Prompt task fixes its cue type and color semantics, while
initializations vary the object or position referred to by the cue.
Cue colors are selected according to scene contrast. The markers modify
only the policy's image inputs and do not alter physical collision
relations.

Before formal comparisons, representative trajectories are manually
reviewed to verify that object identity, process constraints, termination,
and stability clauses are consistent with the task semantics. Physically
invalid initializations or semantically ambiguous instances are repaired or
excluded before the final aggregate is computed.

\subsection{Per-Step Trajectory-Contract Satisfaction}
\label{app:smart-contract}

LIBERO-SMART represents requirements as a trajectory-level contract that
jointly captures target correctness, process consistency, and result
stability. For task \(i\), let \(\tau_{\leq t}\) denote the execution
history through step \(t\). Strict contract satisfaction at that step is
\begin{equation}
\label{eq:smart-contract}
C_i(t)=\mathbf{1}\!\left[G_i(s_t)\land H_i(\tau_{\leq t})\land R_i(\tau_{\leq t})\right],
\end{equation}

where \(G_i\) requires the current result to satisfy the current
instruction, \(H_i\) requires applicable object restrictions,
prohibitions, and step order throughout the observed history, and \(R_i\)
requires a semantically valid and stable result. These terms answer ``is
the result correct?'', ``was the instruction followed?'', and ``was the
task actually completed?'', rather than treating a single transient
target state as complete instruction following.

Task goals and applicable prohibitions are evaluated from environment
state and predicates. Requirements involving execution history are
checked from the trajectory, including forbidden events, state changes
of protected objects, wrong-object manipulation, and the order of
multi-step events. Tasks without a corresponding language constraint do
not receive an additional restriction.

Termination conditions are task-specific. Placement requires release and
stable support; grasping or holding requires the correct grasp and
corresponding stable state; switching tasks check device state and
instruction. Stability checks require applicable state conditions to
hold for the number of consecutive steps specified by the task.
Trajectory tasks use more permissive criteria for trajectory matching,
release, and stability and do not require pointwise replay of the drawn
curve.

The evaluation interface receives actions generated from images, robot
state, and the current instruction. Task binding and contract checking
are performed by an independent evaluator. Additional simulator
information used to check target and trajectory constraints is not
provided as policy input.

A forbidden event remains part of the observed history even if the final
state is later repaired. For example, placing a non-target object in a
basket and removing it still violates an instruction that forbids placing
other objects there. Step-order requirements likewise depend on the
recorded event sequence, not only on the set of completed operations.

\subsection{Interaction Horizon and Success Aggregation}
\label{app:smart-aggregation}

End-to-end evaluation requires attaining strict success within the
original interaction horizon and maintaining the task contract for 30
full subsequent environment action steps, rather than ending the check
at a single successful state. The policy continues to execute actions
conditioned on current observations during the holding check.

Let \(C_i(t)\) denote strict contract satisfaction at step \(t\), as
defined in Equation~(\ref{eq:smart-contract}), and let \(T_i\) denote
the original interaction horizon. Let \(\mathcal{T}_i\) be the set of
environment steps in the task's valid evaluation interval that do not
exceed \(T_i\). The success indicator is
\begin{equation}
\label{eq:smart-success-30}
S_i^{30}=\mathbf{1}\!\left[
\exists\,t\in\mathcal{T}_i:\quad
\prod_{k=0}^{30}C_i(t+k)=1
\right].
\end{equation}

The holding window in Equation~(\ref{eq:smart-success-30}) may extend beyond the original horizon,
but must
start from a successful state at or before \(T_i\) and cover all 30
subsequent steps; evaluation therefore extends at most to \(T_i+30\).
Extra steps only complete a holding check and cannot initiate a new
success window beyond the horizon.

If a candidate window is interrupted without violating history-dependent
constraints, the policy may attain success again and start another
window within the original horizon. Object restrictions, forbidden
events, and step-order history persist throughout the episode and are
not reset by a retry. Success requires every step in a window to satisfy
the same strict contract, including applicable history-dependent
constraints and task-specific stability requirements. If no such window
exists, the episode fails. Ever-success and strict terminal success are
distinct auxiliary statistics, not substitutes for this indicator.

For agentic systems, let \(t_i^{\mathrm{end}}\) be the environment step at
which the budget is exhausted or the upper-level VLM reports completion,
whichever occurs first. The success indicator is
\(S_i^{\mathrm{agent}}=C_i(t_i^{\mathrm{end}})\), evaluated independently
at termination. No additional 30-step window is imposed. The contract's
task-specific stability and release or holding conditions still apply at
termination, and history-dependent constraints are checked over the entire
episode. Neither an earlier successful state nor the VLM's completion
report alone establishes success. All three agentic systems use this same rule.

As an auxiliary metric, Goal-only records whether the current SMART instruction's positive goals are
simultaneously satisfied at any post-action state within the episode's original step budget and
before termination. It preserves the specified objects, targets, and instantaneous spatial
relations, but omits forbidden-event history, action-order constraints, displacement-preservation
guards, release checks, velocity thresholds, and stability or holding windows. It is distinct from
both ever-strict success and success on the original LIBERO task. The column averages this indicator
over the three suites.

Results are first averaged over initialization instances for each task,
then aggregated with equal task weight within each suite and equal suite
weight across the three suites. Within each evaluation protocol, models
use identical tasks, initializations, interaction budgets, and task
contracts; the end-to-end and agentic evaluation settings are specified
in Appendix~\ref{agentic-calling-protocol}.

\subsection{Relationship Between Training Data and Evaluation Tasks}
\label{app:smart-training-boundary}

During LIBERO target-domain adaptation, supervision is derived
from the original LIBERO demonstrations. We perform action-aligned
segmentation and generate local instructions, linguistic rewrites, and
spatial visual supervision from the corresponding observations and
action segments; execution trajectories of LIBERO-SMART evaluation
tasks are not used as training demonstrations. The additional
supervision therefore expands the instruction expressions and visual
conditions associated with existing manipulation experience, rather
than introducing demonstrations collected for the evaluation tasks.

LIBERO-SMART shares LIBERO's basic manipulation skills and scene assets,
but its evaluation instances do not reuse the corresponding training
trajectories. The benchmark changes the conditions under which familiar
skills are invoked through held-out task combinations, language
conditions, object--target relations, and macro scene layouts. Basic
primitives such as grasping, lifting, and placement may therefore be
shared, while their phase-specific combinations and semantic bindings are
new at evaluation time.

The exact evaluation instructions are new instantiations rather than
copies of the original episode captions. They combine object roles,
target relations, ordering, exclusion, and phase-specific termination
requirements that are not present in the corresponding training
episodes. Likewise, evaluation visual cues use varied colors and
locations, and their association with the current object or target is
changed across instances. Circles, points, regions, and trajectories are
thus not tied to a fixed coordinate or a fixed object identity.

This separation preserves the controllability provided by familiar
objects and basic skills while preventing a fixed appearance or
object--location association from being a sufficient solution. The
benchmark consequently measures whether the policy can apply learned
instruction-following patterns to new semantic, relational, and visual
configurations, rather than directly imitating evaluation trajectories.


